\section{Rational points on a circle (Number theory)}
Could there be a circle on the plane where only one point has rational co-ordinates? A point $p=(x,y)$ has rational co-ordinates iff $x,y\in\mathbb{Q}$. What about a sphere? Are there circles with more than one rational point? Are there circles with infinite rational points? Can you find a circle with exactly $n$ rational points on it for each $n\in\mathbb{N}$?

\subsection{Solution}
Yes, there could. Consider the circle: $(x-r)^2+y^2=r^2$ for some irrational $r$.
$(0,0)$ is a point on this circle and is rational. But from the equation it follows
that: $x^2-2xr+r^2+y^2=r^2$ or $x^2+y^2=2xr$ or $r=(x^2+y^2)/2x$. If there
was a rational solution to this equation where $x\ne0$ then it would follow that
$r$ is rational since it is the result of multiplication, division, addition and subtraction
of rational numbers and this is a contradiction. This means this equation can only
have rational solutions for $x=0$. If $x=0$ then $r^2+y^2=r^2$ and so $y=0$ and so
$(0,0)$ is the only such solution.

Same solution applies to a sphere: $(x-r)^2+y^2+z^2=r^2$ for some irrational $r$.
In this case: $r=(x^2+y^2+z^2)/2x$. Which means that for $x\ne0$ any rational solution would imply a rational $r$.
$(0,0,0)$ is therefore the only rational solution.

Yes, there are circles with more than one rational point, and as soon as a circle has three rational
points it has infinitely many. Three points which are not on one line determine a circle. Its center is
where the perpendicular bisectors of two of the chords meet, and since these are lines with rational
coefficients the center $(a,b)$ is rational, and so is $r^2$, the squared distance from the center to one
of the points. Now take one rational point $P=(p,q)$ on the circle and draw through it the line
$y-q=t(x-p)$ for some rational slope $t$. Substituting it into $(x-a)^2+(y-b)^2=r^2$ gives a quadratic
equation in $x$ with rational coefficients. One of its roots is $p$, and the sum of its two roots is
rational (it is minus the ratio of two of the coefficients), so the other root is rational too, and so is
the second point where the line meets the circle. Different slopes give different points, so there are
infinitely many. For the unit circle, starting from $P=(-1,0)$, this gives the points
$\left(\frac{1-t^2}{1+t^2},\frac{2t}{1+t^2}\right)$ for $t\in\mathbb{Q}$, for example $(3/5,4/5)$ for $t=1/2$.
And the circle $x^2+y^2=2$, which has the rational points $(1,1),(-1,1),(1,-1),(-1,-1)$, also has
infinitely many, for example $(7/5,1/5)$.

So a circle has $0$, $1$, $2$ or infinitely many rational points, and for $n\ge3$ there is no circle with
exactly $n$ of them. The other three counts all happen. $1$ was shown above. For $0$ take the circle
$(x-\sqrt{2})^2+y^2=1$: a rational point on it would satisfy $x^2+y^2+1=2\sqrt{2}x$, which for $x\ne0$
would make $\sqrt{2}$ rational, and for $x=0$ gives $y^2=-1$. For $2$ take the circle with center
$(1/2,\sqrt{2})$ through $(0,0)$ and $(1,0)$, that is $(x-1/2)^2+(y-\sqrt{2})^2=9/4$: a rational point on
it would satisfy $x^2-x+y^2=2\sqrt{2}y$, which for $y\ne0$ would make $\sqrt{2}$ rational, and for $y=0$
leaves only $x=0$ and $x=1$.
